\documentclass[10pt]{article}
\usepackage[margin=0.84in]{geometry}
\usepackage[T1]{fontenc}
\usepackage{amsmath,amssymb,amsthm}
\usepackage{graphicx}
\usepackage{enumitem}
\usepackage[colorlinks=true,linkcolor=blue,citecolor=blue,urlcolor=blue]{hyperref}

\newtheorem{theorem}{Theorem}
\newtheorem{lemma}{Lemma}
\newtheorem{corollary}{Corollary}
\newcommand{\T}{\mathbb T}
\newcommand{\Pp}{\mathcal P}
\newcommand{\wh}[1]{\widehat{#1}}

\title{The strict Courtney--Sarason lower bound\\
for every self-adjoint Toeplitz truncation level}
\author{Run \texttt{fa\_banach\_001} --- \texttt{agent\_lane\_11}}
\date{13 August 2026}

\begin{document}
\maketitle

\noindent\textbf{Status.} Candidate full solution of Question~1 in
Courtney--Sarason~\cite{courtney-sarason}; likely valid subject to human review.
Questions~2--4 of the source are not claimed solved.

\section{Source question and notation}

For $f\in L^\infty(\T)$, let
\[
 A_{f,N}=\bigl(\wh f(j-\ell)\bigr)_{j,\ell=0}^{N}
\]
be the compression of multiplication by $f$ to the polynomials
$\Pp_N$. For a self-adjoint Toeplitz matrix $A$, let $c_A$ be the least
$L^\infty$ norm of an inducing function with the prescribed Fourier
coefficients, and set
\[
 c_N=\max_{0\ne A=A^*\text{ Toeplitz}}
          \frac{c_A}{\|A\|}.
\]
The source asks in Section~8, PDF page~13:

\begin{center}
\includegraphics[width=0.96\textwidth]{figures/open_problem_crop.png}
\end{center}

Courtney--Sarason prove that the minimum-norm inducer is an alternating step
function and, conversely, that every alternating step function of height one
and order at most $N$ is the minimum-norm inducer for its compression
$A_{f,N}$~\cite[Sections 4--5]{courtney-sarason}. Consequently
\begin{equation}\label{eq:step-lower}
 c_N\ge \frac1{\|A_{f,N}\|}
\end{equation}
for every such height-one step function $f$.

\section{Proof intuition}

The source dilates a symbol $f$ to $f(z^k)$ only at the exact multiple level
$kN$. The residue-class proof actually loses nothing at an arbitrary level
$M$: each residue modulo $k$ carries a Toeplitz block for $f$, and the largest
block has level $\lfloor M/k\rfloor$. Thus a single strict level-two witness
can be copied without changing its compression norm into every level
$M\ge4$ by taking $k=\lfloor M/2\rfloor$. The only uncovered level is $3$,
where the source already supplies a rigorous strict witness.

\section{The non-multiple dilation lemma}

\begin{lemma}\label{lem:dilation}
Let $f\in L^\infty(\T)$, let $M,k$ be positive integers with $k\le M$, and
put $g(e^{i\theta})=f(e^{ik\theta})$. Then
\begin{equation}\label{eq:dilation}
 \|A_{g,M}\|=\|A_{f,\lfloor M/k\rfloor}\|.
\end{equation}
More precisely, $A_{g,M}$ is unitarily equivalent to the direct sum
\begin{equation}\label{eq:blocksum}
 \bigoplus_{r=0}^{\min\{k-1,M\}}
 A_{f,\lfloor(M-r)/k\rfloor}.
\end{equation}
\end{lemma}

\begin{proof}
The Fourier coefficient $\wh g(n)$ is zero unless $k$ divides $n$, and
$\wh g(kj)=\wh f(j)$. For each residue $r\in\{0,\ldots,k-1\}$, let $S_r$
be the span of the monomials
\[
 z^r,z^{r+k},\ldots,z^{r+m_rk},
 \qquad m_r=\left\lfloor\frac{M-r}{k}\right\rfloor,
\]
omitting empty spaces. The spaces $S_r$ are mutually orthogonal, span
$\Pp_M$, and reduce $A_{g,M}$. Under the unitary identification
$z^{r+jk}\leftrightarrow z^j$, the $(p,q)$ entry of the restriction is
\[
 \wh g(k(p-q))=\wh f(p-q),
\]
so the block is $A_{f,m_r}$. This proves \eqref{eq:blocksum}. The $r=0$
block has level $m_0=\lfloor M/k\rfloor$, while every other block is a
principal compression of it. Therefore the norm of the direct sum is exactly
the norm of the $r=0$ block, proving \eqref{eq:dilation}.
\end{proof}

The exact-multiple case $M=kN$ is Proposition~7.1 of the source. The arbitrary
$M$ formulation gives an additional comparison.

\begin{corollary}\label{cor:comparison}
For positive integers $M,k$ with $k\le M$,
\[
 c_M\ge c_{\lfloor M/k\rfloor}.
\]
\end{corollary}

\begin{proof}
Take an alternating step maximizer $f$ at level $\lfloor M/k\rfloor$ and
put $g(z)=f(z^k)$. Its order is at most
$k\lfloor M/k\rfloor\le M$, so \eqref{eq:step-lower} applies at level $M$.
Lemma~\ref{lem:dilation} preserves the relevant compression norm.
\end{proof}

\section{A rigorous strict seed at level two}

Let $\psi$ be the even alternating step function of height one and order two
which equals $1$ at $1\in\T$ and has discontinuities at the angles
\[
 \pm L,\quad \pi\pm R,
 \qquad L=\frac{\pi}{15},\quad R=\frac{\pi}{3}.
\]
Direct integration (the same formula used in Section~7 of the source) gives
\begin{align}
 a_0:=\wh\psi(0)&=-\frac15,\label{eq:a0}\\
 a_1:=\wh\psi(1)&=\frac2\pi
   \left(\sin\frac{\pi}{15}-\frac{\sqrt3}{2}\right),\label{eq:a1}\\
 a_2:=\wh\psi(2)&=\frac1\pi
   \left(\frac{\sqrt3}{2}+\sin\frac{2\pi}{15}\right).\label{eq:a2}
\end{align}
Thus
\[
 A_{\psi,2}=
 \begin{pmatrix}a_0&a_1&a_2\\a_1&a_0&a_1\\a_2&a_1&a_0\end{pmatrix}.
\]
Put $x=|a_1|$ and $y=a_2$. The antisymmetric eigenvector gives the eigenvalue
$-1/5-y$, and the two symmetric eigenvalues are
\begin{equation}\label{eq:eigenvalues}
 -\frac15+\frac y2\ \pm\ \frac12\sqrt{y^2+8x^2}.
\end{equation}

For completeness, all numerical separation needed below follows from the
rational bounds
\begin{gather*}
 3.14159<\pi<3.14160,\qquad
 0.20791<\sin(\pi/15)<0.20792,\\
 0.40673<\sin(2\pi/15)<0.40674,\qquad
 0.86602<\sqrt3/2<0.86603.
\end{gather*}
Machin's formula and alternating series prove the first bound; alternating
Taylor bounds prove the two sine bounds; the last follows by squaring. They
imply
\[
 x<0.41898,\qquad 0.40512<y<0.40514.
\]
Hence the antisymmetric eigenvalue has modulus below $0.60515$, while
\[
 \left|-\frac15+\frac y2\right|<0.00257,
 \qquad
 \frac12\sqrt{y^2+8x^2}<0.62620.
\]
By \eqref{eq:eigenvalues},
\begin{equation}\label{eq:strict-seed}
 \|A_{\psi,2}\|<0.62877<\frac2\pi.
\end{equation}
The accompanying verifier certifies each displayed decimal inequality by
exact rational arithmetic. From \eqref{eq:step-lower} and
\eqref{eq:strict-seed}, $c_2>\pi/2$.

\section{Full answer to Question 1}

\begin{theorem}\label{thm:main}
For every integer $N>1$, the Courtney--Sarason mini-max constant satisfies
\[
 c_N>\frac\pi2.
\]
\end{theorem}

\begin{proof}
The case $N=2$ is \eqref{eq:strict-seed}. For $N=3$, Section~7 of the source
constructs an alternating step function proving
\[
 c_3\ge\frac{\pi\kappa}{2\sqrt2},
\]
where $\kappa$ is the largest root of
$1-3x-3x^2+3x^3$. Since this polynomial is negative at $3/2$,
$\kappa>3/2$; consequently
$c_3>3\pi/(4\sqrt2)>\pi/2$.

Now let $N\ge4$ and set $k=\lfloor N/2\rfloor$. Then
\[
 \left\lfloor\frac Nk\right\rfloor=2,
 \qquad 2k\le N.
\]
Equivalently, the dilated function $g(z)=\psi(z^k)$ is an alternating step
function of order $2k\le N$, and Lemma~\ref{lem:dilation} gives
\[
 \|A_{g,N}\|=\|A_{\psi,2}\|<\frac2\pi.
\]
Applying \eqref{eq:step-lower} to $g$ yields $c_N>\pi/2$.
\end{proof}

\section{Verification, novelty, and limitations}

The standard-library verifier proves the rational interval certificate, then
checks the residue block levels for every $4\le N\le1000$. The finite sweep is
only a sanity check; Lemma~\ref{lem:dilation} proves the statement for all
$N$.

A bounded run-index and arXiv/web search on 13 August 2026 used the exact title,
authors, $c_N$, ``self-adjoint Toeplitz $\pi/2$,'' ``Courtney--Sarason,'' and
the exact wording of Questions~1--2. It recovered the source and later papers
citing it for other Toeplitz/Blaschke questions, but no resolution or
non-multiple dilation formula. The novelty assessment is therefore
\emph{apparently new within the bounded search}, not an exhaustive guarantee.

Theorem~\ref{thm:main} completely answers Question~1. Corollary
\ref{cor:comparison} gives new non-adjacent comparisons, but it does not imply
$c_{N+1}>c_N$. Questions~2--4 remain open. Human review should prioritize the
source-derived extremal reformulation and the arbitrary-level block
decomposition.

\begin{thebibliography}{9}
\bibitem{courtney-sarason}
D. Courtney and D. Sarason,
\emph{A mini-max problem for self-adjoint Toeplitz matrices},
Math. Scand. \textbf{110} (2012), 82--98;
\href{https://doi.org/10.7146/math.scand.a-15198}{doi:10.7146/math.scand.a-15198};
\href{https://arxiv.org/abs/1009.2553}{arXiv:1009.2553}.
\end{thebibliography}

\end{document}
